% Part: first-order-logic
% Chapter: tableaux
% Section: proof-theoretic-notions

\documentclass[../../../include/open-logic-section]{subfiles}

\begin{document}

\iftag{FOL}
      {\olfileid{fol}{tab}{ptn}}
      {\olfileid{pl}{tab}{ptn}}
      
\olsection{Gagasan Teoretis-Bukti}

\begin{editorial}
Bagean iki nglumpukake definisi relasi provabilitas lan konsistensi
kanggo tableaux.
\end{editorial}

\begin{explain}
Kaya nalika kita nemtokake sawetara gagasan semantis wigati
(validitas, akibat, lan bisa dicukupi), saiki kita nemtokake
\emph{gagasan teoretis-bukti} kang cocog. Gagasan iki ora ditemtokake
kanthi nggunakake panyukupan !!{sentence}s ing !!{structure}s, nanging
kanthi nggunakake anane !!{tableau} katutup tartamtu. Panemuan yen
gagasan-gagasan kasebut padha iku panemuan wigati. Kasunyatan manawa
gagasan kasebut padha dadi isi \emph{teorema kasahihan} lan
\emph{teorema kelengkapan}.
\end{explain}


\begin{defn}[Teorema]
!!{sentence}~$!A$ iku \emph{teorema} yen ana !!{tableau} katutup
kanggo~$\sFmla{\False}{!A}$. Kita nulis $\Proves !A$ yen $!A$
iku teorema, lan $\Proves/ !A$ yen dudu teorema.
\end{defn}

\begin{defn}[!!^{derivability}]
!!^a{sentence} $!A$ \emph{bisa !!{derivable} saka} himpunan
!!{sentence}s~$\Gamma$, ditulis $\Gamma \Proves !A$, yen lan mung yen
ana himpunan winates $\{!B_1, \dots, !B_n\} \subseteq \Gamma$
lan !!{tableau} katutup kanggo himpunan
\[
\{
  \sFmla{\False}{!A},
  \sFmla{\True}{!B_1}, \dots,
  \sFmla{\True}{!B_n}
  \}.
\]
Yen $!A$ ora bisa !!{derivable} saka $\Gamma$, kita nulis $\Gamma
\Proves/ !A$.
\end{defn}

\begin{defn}[Konsistensi]
Himpunan !!{sentence}s~$\Gamma$ iku \emph{ora konsisten} yen lan mung
yen ana himpunan winates $\{!B_1, \dots, !B_n\} \subseteq \Gamma$
lan !!{tableau} katutup kanggo himpunan
\[
\{
  \sFmla{\True}{!B_1}, \dots,
  \sFmla{\True}{!B_n}  
  \}.
\]
Yen $\Gamma$ dudu himpunan kang ora konsisten, kita ngarani himpunan iku
\emph{konsisten}.
\end{defn}

\begin{prop}[Refleksivitas]
\ollabel{prop:reflexivity}
Yen $!A \in \Gamma$, mula $\Gamma \Proves !A$.
\end{prop}

\begin{proof}
Yen $!A \in \Gamma$, $\{!A\}$ iku subhimpunan winates saka~$\Gamma$,
lan !!{tableau}
\begin{oltableau}
  [\sFmla{\False}{\formula{A}}, just = \TAss
    [\sFmla{\True}{\formula{A}}, just = \TAss,close]
  ]
\end{oltableau}
iku katutup.
\end{proof}
  

\begin{prop}[Monotonisitas]
\ollabel{prop:monotonicity}
Yen $\Gamma \subseteq \Delta$ lan $\Gamma \Proves !A$, mula $\Delta
\Proves !A$.
\end{prop}

\begin{proof}
Saben subhimpunan winates saka~$\Gamma$ uga dadi subhimpunan winates
saka~$\Delta$.
\end{proof}

\begin{prop}[Transitivitas]
\ollabel{prop:transitivity}
Yen $\Gamma \Proves !A$ lan $\{!A\} \cup
\Delta \Proves !B$, mula $\Gamma \cup \Delta \Proves !B$.
\end{prop}

\begin{proof}
Yen $\{!A\} \cup \Delta \Proves !B$, mula ana subhimpunan winates
$\Delta_0 = \{!C_1, \dots, !C_n\} \subseteq \Delta$ saengga
\begin{align*}
\{\sFmla{\False}{!B}, & \sFmla{\True}{!A}, \sFmla{\True}{!C_1},
\dots, \sFmla{\True}{!C_n}\}
\intertext{nduweni !!{tableau} katutup. Yen $\Gamma \Proves !A$,
  mula ana $!D_1$, \dots, $!D_m \in \Gamma$ saengga}
\{\sFmla{\False}{!A}, & \sFmla{\True}{!D_1},
\dots, \sFmla{\True}{!D_m}\}
\end{align*}
nduweni !!{tableau} katutup. Cathetan koreksi sumber OLPL-016: sumber
Inggris menehi tandha subhimpunan marang dhaptar ukara tanpa kurung
kurawal; definisi sadurunge mbutuhake saben ukara kasebut kalebu ing
himpunan premis, kaya kang ditulis ing kene.

Saiki gatekna !!{tableau} kanthi asumsi
\[
\sFmla{\False}{!B},
\sFmla{\True}{!C_1}, \dots, \sFmla{\True}{!C_n},
\sFmla{\True}{!D_1}, \dots, \sFmla{\True}{!D_m}.
\]
Terapna aturan \Cut{} marang~$!A$. Iki ngasilake rong cabang: siji
ngemot $\sFmla{\True}{!A}$, lan sijine ngemot
$\sFmla{\False}{!A}$. Mula, ing cabang kapisan, kabeh
\[
\{\sFmla{\False}{!B}, \sFmla{\True}{!A},
\sFmla{\True}{!C_1}, \dots, \sFmla{\True}{!C_n}\}
\]
wis cumawis. Amarga ana !!{tableau} katutup kanggo asumsi-asumsi iki,
kita bisa nggandhengake tableau iku ing cabang kasebut; saben cabang
kang ngliwati $\sFmla{\True}{!A}$ bakal katutup. Ing cabang liyane,
kabeh
\[
\{\sFmla{\False}{!A}, \sFmla{\True}{!D_1}, \dots,
\sFmla{\True}{!D_m}\}
\]
wis cumawis, mula kita uga bisa njangkepi sisih liyane kanggo entuk
!!{tableau} katutup. Iki nuduhake $\Gamma \cup \Delta \Proves !B$.
\end{proof}

Gatekna yen iki mligine ateges: yen $\Gamma \Proves !A$ lan $!A
\Proves !B$, mula $\Gamma \Proves !B$. Uga ndherek yen $!A_1,
\dots, !A_n \Proves !B$ lan $\Gamma \Proves !A_i$ kanggo saben~$i$,
mula $\Gamma \Proves !B$.

\begin{prop}
\ollabel{prop:incons}
$\Gamma$ ora konsisten yen lan mung yen $\Gamma \Proves !A$ kanggo
saben !!{sentence}~$!A$.
\end{prop}

\begin{proof}
Gladhen.
\end{proof}

\tagprob{FOL}
\begin{prob}
Buktekna \olref[fol][tab][ptn]{prop:incons}
\end{prob}
\tagendprob

\tagprob{notFOL}
\begin{prob}
Buktekna \olref[pl][tab][ptn]{prop:incons}
\end{prob}
\tagendprob

\begin{prop}[Kekompakan]
  \ollabel{prop:proves-compact}
  \begin{enumerate}
  \item Yen $\Gamma \Proves !A$, mula ana subhimpunan winates
    $\Gamma_0 \subseteq \Gamma$ saengga $\Gamma_0 \Proves !A$.
  \item Yen saben subhimpunan winates saka~$\Gamma$ konsisten,
    mula $\Gamma$ konsisten.
  \end{enumerate}
\end{prop}

\begin{proof}
  \begin{enumerate}
    \item Yen $\Gamma \Proves !A$, mula ana subhimpunan winates
      $\Gamma_0 = \{!B_1, \dots, !B_n\}$ lan !!{tableau} katutup
      kanggo
      \[
      \{\sFmla{\False}{!A}, \sFmla{\True}{!B_1}, \dots, \sFmla{\True}{!B_n}\}
      \]
      !!{tableau} iki uga nuduhake $\Gamma_0 \Proves !A$.
    \item Yen $\Gamma$ ora konsisten, mula kanggo sawijining
      subhimpunan winates $\Gamma_0 = \{!B_1, \dots, !B_n\}$ ana
      !!{tableau} katutup kanggo
      \[
      \{\sFmla{\True}{!B_1}, \dots, \sFmla{\True}{!B_n}\}
      \]
      !!{tableau} katutup iki nuduhake yen $\Gamma_0$ ora konsisten.
  \end{enumerate}
\end{proof}

\end{document}
